\section{提示词工程：指令的质量决定产出的质量}

\subsection{好指令的三要素}

\begin{importantbox}{目标 + 约束 + 上下文}
\begin{itemize}[leftmargin=1.5em]
  \item \textbf{目标}：明确要做什么（"在表格上方加搜索框"）
  \item \textbf{约束}：限定实现方式（"走后端 API，debounce 300ms"）
  \item \textbf{上下文}：引用具体文件、样式规范、现有组件
\end{itemize}
差的指令："给用户列表加个搜索功能"。好的指令减少了歧义，Claude 不需要猜你要前端还是后端搜索。
\end{importantbox}

\subsection{四个官方推荐技巧}

\begin{enumerate}[leftmargin=1.5em]
  \item \textbf{用 \texttt{@} 引用具体文件}——\texttt{@src/auth/login.ts} 比"那个认证相关的代码"精确得多，还可以指定行号范围 \texttt{@src/auth/login.ts\#5-30}。
  \item \textbf{贴截图和图片}——调试 UI 问题时，\texttt{Ctrl+V} 粘贴"现在长这样"和"我想要这样"的截图，比文字描述有效十倍。
  \item \textbf{给验证标准}——给 Claude 一个可以自我验证的标准（测试用例、截图、预期输出），它就能自己检查工作质量，减少来回修改。
  \item \textbf{明确"不要做"的事}——Claude 有时会过度热心，你让它修 bug，它顺手把周围代码也重构了。明确说"只改这个函数，不要动其他代码"。
\end{enumerate}

\begin{knowledgebox}{模糊指令也有用武之地}
探索性任务用模糊指令反而更好，比如"这个文件有什么可以改进的地方？"、"帮我理解一下这个模块的架构"。但到了实现阶段，要切回精确模式。
\end{knowledgebox}

\subsection{本章小结}

好的提示词 = 目标 + 约束 + 上下文 + 验证标准。探索时用模糊指令，实现时用精确指令。

%% ====================================================================
